% Compile from the repository root: xelatex -output-directory=build examples/model-comparison.tex
\documentclass[11pt]{article}
\usepackage[paperwidth=180mm,paperheight=155mm,margin=16mm]{geometry}
\usepackage[T1]{fontenc}
\usepackage{lmodern,booktabs,array}
\usepackage{gradbars}
\pagestyle{empty}
\renewcommand{\familydefault}{\sfdefault}
\setlength{\parindent}{0pt}
\begin{document}
{\small\color{gradbarsBlue}\bfseries GRADBARS / 01 \quad MODEL COMPARISON}\par
\vspace{5mm}
{\LARGE\bfseries Numbers, with perspective.}\par
\vspace{3mm}
{\color{gradbarsInk}A compact visual layer for research tables.}\par
\vspace{8mm}
\gradbarssetup{width=27mm,height=1.5ex,label width=3.6em,precision=1}
\renewcommand{\arraystretch}{1.65}
\begin{tabular*}{\linewidth}{@{\extracolsep{\fill}}lcc@{}}
\toprule
\textbf{Model} & \textbf{Accuracy (\%)} & \textbf{Latency (ms)}\\
\midrule
Baseline     & \gradbar{72.4} & \gradbar[theme=teal,max=50]{18.2}\\
Compact      & \gradbar{81.6} & \gradbar[theme=teal,max=50]{12.5}\\
Transformer  & \gradbar{89.3} & \gradbar[theme=teal,max=50]{42.8}\\
\textbf{Proposed} & \gradbar{92.7} & \gradbar[theme=teal,max=50]{24.1}\\
\bottomrule
\end{tabular*}\par
\vspace{5mm}
{\footnotesize\color{gradbarsInk}
Illustrative data. Accuracy: higher is better; latency: lower is better.
Every bar starts at zero. Each column shares one maximum: 100 and 50.}\par
\vspace{7mm}
{\small\bfseries One command per cell.}\par
\vspace{2mm}
{\small\ttfamily\string\gradbar\{92.7\}\quad
\string\gradbar[max=50,theme=teal]\{24.1\}}
\end{document}
